\chapter{Evaluation and Results}

\section{Evaluation Setup}
A comparative usability study measured educator-oriented authoring tasks under two conditions: AI-ActivEdu and a WordPress-based H5P workflow. Participants completed equivalent scenarios under controlled conditions.

\section{Quantitative Results}
The measured outcomes indicate substantial gains for the proposed platform:
\begin{itemize}[leftmargin=2em]
\item Task completion: 95\% versus 53.3\% baseline.
\item SUS score: 84.3 versus 36.7.
\item Average authoring time: 3.4 minutes versus 46.4 minutes.
\item Higher first-attempt correctness and lower recovery effort.
\end{itemize}

% Comparative evaluation metrics — four bar charts
% pgfplots groupplot layout (2x2)
\begin{tikzpicture}
\begin{groupplot}[
  group style={
    group size=2 by 2,
    horizontal sep=2.2cm,
    vertical sep=2.8cm,
  },
  ybar,
  bar width=16pt,
  width=6.8cm, height=5.5cm,
  symbolic x coords={New Platform, WordPress},
  xtick=data,
  xticklabel style={font=\scriptsize, rotate=18, anchor=east},
  nodes near coords,
  nodes near coords style={font=\scriptsize},
  enlarge x limits=0.55,
  axis lines*=left,
  ymajorgrids=true,
  grid style={dashed, gray!40},
]

%% Plot (a): Task Completion Rate
\nextgroupplot[
  title={\small\textbf{(a) Task Completion Rate (\%)}},
  ymin=0, ymax=110,
  ytick={0,20,40,60,80,100},
  ylabel={\scriptsize Completion (\%)},
]
\addplot[fill=blue!45, draw=blue!70!black] coordinates
  {(New Platform, 95) (WordPress, 30)};

%% Plot (b): Median Task Time
\nextgroupplot[
  title={\small\textbf{(b) Median Task Time (min)}},
  ymin=0, ymax=55,
  ytick={0,10,20,30,40,50},
  ylabel={\scriptsize Time (minutes)},
]
\addplot[fill=green!45, draw=green!70!black] coordinates
  {(New Platform, 3.2) (WordPress, 45)};

%% Plot (c): Mean SUS Score
\nextgroupplot[
  title={\small\textbf{(c) Mean SUS Score (0--100)}},
  ymin=0, ymax=110,
  ytick={0,20,40,60,80,100},
  ylabel={\scriptsize SUS Score},
  extra y ticks={68},
  extra y tick style={
    tick label style={font=\tiny, color=red!70},
    grid=major,
    grid style={dashed, red!40},
  },
  extra y tick labels={avg\,(68)},
]
\addplot[fill=orange!45, draw=orange!70!black] coordinates
  {(New Platform, 82.5) (WordPress, 38.0)};

%% Plot (d): User Satisfaction
\nextgroupplot[
  title={\small\textbf{(d) User Satisfaction (1--5)}},
  ymin=0, ymax=6,
  ytick={0,1,2,3,4,5},
  ylabel={\scriptsize Satisfaction score},
]
\addplot[fill=purple!45, draw=purple!70!black] coordinates
  {(New Platform, 4.8) (WordPress, 2.1)};

\end{groupplot}
\end{tikzpicture}


\section{Qualitative Findings}
Post-task interviews identified recurring themes:
\begin{enumerate}[leftmargin=2em]
\item The timeline model aligned with educator mental models of teaching sequences.
\item Immediate preview reduced uncertainty and repeated trial-and-error cycles.
\item AI suggestions accelerated drafting but required teacher review for precision.
\item Participants preferred explicit workflow guidance to plugin-style configuration menus.
\end{enumerate}

\section{Statistical and Practical Interpretation}
Even with modest sample size, effect magnitudes were large enough to indicate practical significance for institutional operations. Time savings scale materially when multiplied by weekly content production loads across departments.

\section{Threats to Evaluation Validity}
Key limitations include sample size and context constraints. Results should be interpreted as strong directional evidence, with broader deployments recommended for multi-institution confirmation.
